\documentclass[12pt]{article}
\usepackage[T1]{fontenc}
\usepackage[utf8]{inputenc}
\usepackage{lmodern}
\usepackage{textcomp}
\usepackage{enumitem}
\usepackage{geometry}
\geometry{margin=1in}
\begin{document}
\begin{center}
Answer Key - Economics - K-12 Level Assessment 24 \
\small (Answers are ordered exactly as in the question paper.)
\end{center}

\section*{Multiple Choice}
\begin{enumerate}[label=\arabic*.]
\item \textbf{Answer:} A
\\ \textit{Options:} A) result in higher domestic prices.; B) promote trade between nations.; C) do not necessarily affect domestic prices.; D) affect domestic prices: the former raises them while the latter lowers them.
\\ \textit{Note:} Tariffs and quotas limit the supply of imported goods, which leads to reduced competition and allows domestic producers to raise their prices.
\item \textbf{Answer:} A
\\ \textit{Options:} A) divide nominal GDP by the GDP deflator.; B) multiply real GDP by the GDP deflator.; C) divide real GDP by the GDP deflator.; D) multiply nominal GDP by the GDP deflator.
\\ \textit{Note:} Dividing nominal GDP by the GDP deflator adjusts for inflation, allowing for the calculation of real GDP, which reflects the true value of goods and services without the effects of price changes.
\item \textbf{Answer:} A
\\ \textit{Options:} A) eliminate a recessionary gap; B) reduce inflation; C) reduce the interest rate; D) eliminate an inflationary gap
\\ \textit{Note:} Expansionary fiscal policy involves increasing government spending or decreasing taxes to stimulate economic activity, making it an effective tool to eliminate a recessionary gap by boosting aggregate demand and promoting growth during periods of economic downturn.
\item \textbf{Answer:} A
\\ \textit{Options:} A) Diminishing marginal productivity.; B) Diminishing marginal utility.; C) Increasing marginal utility.; D) Increasing marginal productivity.
\\ \textit{Note:} The correct answer is "diminishing marginal productivity" because it describes the phenomenon where adding more of a variable input, while keeping other inputs constant, leads to progressively smaller increases in output.
\item \textbf{Answer:} C
\\ \textit{Options:} A) NDP will be greater than GDP.; B) NI will be greater than GDP.; C) PI will be greater than NI.; D) DPI will be greater than PI.
\\ \textit{Note:} The correct answer is based on the fact that personal income (PI) includes transfer payments, which can exceed net income (NI) when such payments surpass the combined total of Social Security contributions, corporate taxes, and retained earnings, leading to a higher PI.
\end{enumerate}

\section*{True / False}
\begin{enumerate}[label=\arabic*.]
\setcounter{enumi}{5}
\item \textbf{Answer:} True
\\ \textit{Options:} A) True; B) False
\\ \textit{Note:} Credit cards are not included in the M$_{1}$ money supply because M$_{1}$ consists of liquid monetary assets like cash and checking account balances, while credit cards represent borrowed money rather than actual money in circulation.
\item \textbf{Answer:} False
\\ \textit{Options:} A) True; B) False
\\ \textit{Note:} The statement is false because expansionary fiscal policy can still lead to crowding out if the investment curve is inelastic, meaning that increased government spending may raise interest rates and reduce private investment, regardless of the accompanying monetary policy.
\item \textbf{Answer:} False
\\ \textit{Options:} A) True; B) False
\\ \textit{Note:} The statement is false because the natural rate of unemployment refers to the level of unemployment that exists when the economy is at full employment, which can occur alongside a positive inflation rate due to factors like demand-pull inflation or cost-push inflation.
\item \textbf{Answer:} False
\\ \textit{Options:} A) True; B) False
\\ \textit{Note:} An increase in interest rates generally leads to higher borrowing costs, which can reduce consumer spending and business investment, ultimately shifting the aggregate demand curve to the left rather than to the right.
\item \textbf{Answer:} True
\\ \textit{Options:} A) True; B) False
\\ \textit{Note:} In a monopoly, to sell additional units, the firm must lower the price not only for the additional units but for all previous units sold, resulting in marginal revenue being less than the price indicated by the demand curve.
\end{enumerate}

\section*{Long Form Response}
\begin{enumerate}[label=\arabic*.]
\setcounter{enumi}{10}
\item \textbf{Answer:} Excess capacity in monopolistic competition refers to a situation where firms produce below their optimal output level, meaning they have the ability to produce more without increasing their costs significantly. This characteristic arises because firms in monopolistic competition differentiate their products, leading to a downward-sloping demand curve for each firm. As a result, they often set prices above marginal costs to maximize profits, which can lead to inefficiencies in the market. This pricing strategy means that consumers might pay higher prices for less quantity than they would in a perfectly competitive market, ultimately resulting in a loss of economic efficiency. Thus, while firms can enjoy some market power, the presence of excess capacity can hinder overall societal welfare.
\\ \textit{Note:} correct because it accurately describes excess capacity as the condition where firms in monopolistic competition operate below optimal output levels due to product differentiation, which leads to higher prices than marginal costs and results in market inefficiencies.
\item \textbf{Answer:} When an American consumer purchases an entertainment system manufactured in China, it has a direct impact on U.S. national income accounts. This transaction leads to a decrease in net exports, as net exports are calculated by subtracting the value of imports from exports. Since the consumer is importing the entertainment system, imports rise while exports remain unchanged, resulting in lower net exports. Additionally, this purchase contributes to a fall in the Gross Domestic Product (GDP) of the U.S., as GDP measures the total value of goods and services produced within a country. Overall, both net exports and GDP decline due to this import from China.
\\ \textit{Note:} correct because the purchase of an imported entertainment system increases imports, thereby reducing net exports and negatively impacting the calculation of GDP, which considers net exports as a component.
\end{enumerate}

\vfill
\noindent\textit{End of Answer Key}
\end{document}